% capitulo-3.tex
% Capítulo 3 — Metodología
%
% This chapter describes the planned experimental methodology. Citations
% use the citation keys of the project Zotero library (exported to
% references.bib). Items that depend on decisions not yet finalized
% (hardware details, software versions, model selection, resolutions,
% durations, repetitions) are marked with visible \todo{...} placeholders
% to be completed once determined; % TODO comments mark citation gaps. No
% experimental results are reported in this chapter.

\chapter{Metodología}
\label{chap:methodology}

\section{Plataforma experimental}
\label{sec:platform}

La plataforma experimental es el kit de desarrollo NVIDIA Jetson Orin Nano
Super (P3766), que consta de un módulo Jetson Orin Nano de 8~GB con
disipador montado sobre la placa portadora de referencia
\cite{nvidiacorporation2024}. La designación Super identifica la
configuración de software del kit y no un módulo de hardware distinto: el
mismo módulo de 8~GB admite la configuración original y la Super, y esta
última se habilita mediante el flasheo correspondiente
\cite{nvidiacorporation2024}. La elección del kit responde a su uso como
módulo embebido representativo para la visión por computador en el borde
\cite{scalcon2024, alqahtani2026}. En el resto de esta memoria se emplea la
denominación abreviada «Jetson Orin Nano», que es la que aparece en el
título y en el alcance del trabajo.

La \Cref{tab:platform-configuration} reúne las características del equipo y
las condiciones que se mantienen constantes en todas las corridas.

\begin{table}[htbp]
  \centering
  \caption[Configuración de la plataforma experimental]{Configuración de la
  plataforma experimental. Todos los valores se mantienen constantes a lo
  largo de la campaña de mediciones.}
  \label{tab:platform-configuration}
  \begin{tabular}{@{}p{0.26\linewidth}p{0.66\linewidth}@{}}
    \toprule
    \textbf{Elemento} & \textbf{Configuración fija} \\
    \midrule
    Plataforma & Kit de desarrollo NVIDIA Jetson Orin Nano Super (P3766)
                 \cite{nvidiacorporation2024} \\
    Módulo & Jetson Orin Nano de 8~GB sobre la placa portadora de
             referencia \\
    CPU & Arm Cortex-A78AE v8.2 de seis núcleos \\
    GPU & Ampere con 1024 núcleos CUDA y 32 Tensor Cores \\
    Memoria & 8~GB de LPDDR5 con bus de 128~bits, unificada entre CPU y
              GPU, con un ancho de banda de 102~GB/s \\
    Almacenamiento & SSD NVMe FORESEE XP1000F256G de 256~GB (238,5~GiB
                     utilizables) en la ranura M.2 Key M, que aloja el
                     sistema operativo, el entorno de software, los datos
                     de entrada y los registros. No se utiliza tarjeta SD \\
    Modo de energía & 25~W, fijo en todas las corridas \\
    Enfriamiento & Disipador y ventilador originales del kit, con la curva
                   de control del ventilador por defecto de la plataforma;
                   sin modificación manual de la velocidad ni
                   refrigeración externa \\
    Entrada & Archivo de video almacenado en el NVMe y decodificado
              localmente; no se conecta ninguna cámara \\
    Monitoreo & Temperatura del SoC registrada en todas las corridas, con
                el procedimiento de la \Cref{sec:metric-measurement} \\
    \bottomrule
  \end{tabular}
\end{table}

El modo de energía se fija en 25~W. El rango de potencia configurable del
módulo va de 7~W a 25~W \cite{nvidiacorporation2024} y no es un detalle
accesorio: las evaluaciones publicadas de la familia Orin distinguen
explícitamente el modo de energía al informar tasas de fotogramas
\cite{scalcon2024}, de modo que el modo forma parte de la condición
experimental y se documenta junto con ella. Se fija el extremo superior del
rango porque la carga de inferencia es sostenida y se busca caracterizar la
plataforma en la capacidad de operación declarada más alta; los modos de
menor presupuesto no se emplean porque el objeto de estudio es la
arquitectura de la aplicación y no el efecto del modo de energía. Antes de
cada bloque de mediciones se verifica el modo activo con el comando
\texttt{nvpmodel -q}. Dado que los identificadores numéricos de modo
dependen de la configuración de flasheo, la verificación se hace sobre el
nombre del modo y su presupuesto de potencia, y no sobre el número.

El enfriamiento se mantiene tal como el kit se entrega: disipador y
ventilador originales, con la curva de control del ventilador por defecto de
la plataforma, sin modificación manual de la velocidad ni refrigeración
externa. Bajo carga de inferencia sostenida la temperatura del dispositivo
aumenta de manera medible \cite{tang2026}, de modo que la temperatura del
SoC se registra durante todas las corridas y se contrasta con los puntos de
disparo térmico declarados por la plataforma, lo que permite afirmar si
alguna corrida entró en limitación térmica; el procedimiento de registro se
detalla en la \Cref{sec:metric-measurement}. El laboratorio no dispone de
control climático, por lo que la temperatura ambiente se registra al inicio
y al final de cada sesión de mediciones.

La fuente de entrada es un archivo de video almacenado en el NVMe y
decodificado localmente; no se conecta ninguna cámara. La elección es
deliberada: un archivo entrega la misma secuencia de fotogramas en cada
corrida y en cada configuración, lo que permite comparar las salidas de las
configuraciones A, B y C fotograma a fotograma y elimina la variabilidad de
exposición, ganancia y tasa de cuadro propia de un sensor en vivo. El
archivo se fija en la \Cref{sec:pipeline}; su resolución debe ser al menos
la mayor de las resoluciones experimentales de la
\Cref{sec:configurations}, y la forma en que se sirven los fotogramas se
especifica en la \Cref{sec:procedure}.

Todos los experimentos se ejecutan en este único dispositivo, con los nodos
de las configuraciones de ROS~2 ejecutándose como procesos separados del
sistema operativo en el mismo anfitrión. Ejecutar todos los procesos en un
dispositivo aísla el efecto de la propia arquitectura de ROS~2 de los
efectos de red, que están fuera del alcance de esta memoria.

\section{Entorno de software}
\label{sec:software-environment}

El entorno de software se fija en una sola combinación de versiones, se
mantiene constante en todas las corridas y se verifica con los comandos
indicados en la \Cref{tab:software-environment}. La base es JetPack 7.2.1,
que proporciona una imagen raíz Ubuntu 24.04 LTS (Noble) para \texttt{aarch64}
\cite{nvidiacorporation2026jetpack}. Sobre esa base se instala ROS~2 Jazzy
Jalisco, la distribución con soporte de largo plazo que corresponde a Ubuntu
24.04 y que ofrece paquetes binarios para \texttt{arm64} en esa versión
\cite{rep2000, rosjazzydocs}; de este modo la distribución no se compila desde
el código fuente.

\begin{table}[htbp]
  \centering
  \caption[Configuración del entorno de software]{Configuración del entorno de
  software. Todos los valores se mantienen constantes en las tres
  configuraciones experimentales y se verifican con el comando indicado.}
  \label{tab:software-environment}
  \begin{tabular}{@{}L{0.17\linewidth}L{0.47\linewidth}L{0.30\linewidth}@{}}
    \toprule
    \textbf{Elemento} & \textbf{Valor fijo} & \textbf{Verificación} \\
    \midrule
    Sistema operativo & Ubuntu 24.04 LTS (Noble), \texttt{aarch64}, imagen
      raíz de JetPack 7.2.1 & \texttt{cat /etc/os-release} \\
    JetPack & 7.2.1 (\texttt{nvidia-jetpack} 7.2.1-b49) sobre Jetson Linux
      39.2.1, kernel 6.8 & \texttt{dpkg-query --show nvidia-l4t-core} \\
    ROS~2 & Jazzy Jalisco, paquetes binarios \texttt{arm64} &
      \texttt{printenv ROS\_DISTRO} \\
    Middleware & Fast DDS 2.14.7 seleccionado mediante
      \texttt{rmw\_fastrtps\_cpp} & \texttt{printenv RMW\_IMPLEMENTATION} \\
    Perfil de QoS & uno solo, idéntico en todos los tópicos
      (\Cref{sec:qos}) & código de los nodos \\
    Ejecutor & \texttt{SingleThreadedExecutor} de \texttt{rclpy} en todos los
      nodos (\Cref{sec:ros2-implementation}) & código de los nodos \\
    Intérprete de Python & CPython 3.12.3, el del sistema &
      \texttt{python3 --version} \\
    Bibliotecas de imagen & NumPy 1.26.4 y OpenCV 4.6.0, instalados desde
      Ubuntu 24.04 (\texttt{python3-numpy}, \texttt{python3-opencv}) &
      \texttt{python3 -c 'import cv2'} \\
    Runtime de inferencia & el de la \Cref{sec:pipeline}, idéntico en A, B y
      C & \Cref{sec:pipeline} \\
    Monitoreo & \texttt{tegrastats} para frecuencias, temperatura y consumo &
      \Cref{sec:metric-measurement} \\
    \bottomrule
  \end{tabular}
\end{table}

La instalación sigue un procedimiento único, que se ejecuta una vez y se
registra junto con el artefacto experimental: (i)~flasheo de la imagen ISO de
JetPack 7.2.1 con la configuración Super sobre el NVMe; (ii)~alta del
repositorio oficial de ROS~2 mediante el paquete \texttt{ros2-apt-source} e
instalación de \texttt{ros-jazzy-ros-base} junto con los paquetes que la
cadena requiera, más las bibliotecas de imagen del intérprete
(\texttt{python3-numpy} y \texttt{python3-opencv}); (iii)~congelación de las
versiones con \texttt{apt-mark hold} sobre los paquetes \texttt{ros-jazzy-*}
y los de las bibliotecas de imagen, y registro del manifiesto completo de
paquetes instalados, de modo que una actualización no altere el entorno en
mitad de la campaña; (iv)~selección del middleware con la variable de entorno
\texttt{RMW\_IMPLEMENTATION}; y (v)~verificación con los comandos de la
\Cref{tab:software-environment}.

NumPy y OpenCV se instalan desde los paquetes de Ubuntu y no desde
\texttt{pip}, porque los paquetes del sistema quedan fijados por
\texttt{apt-mark hold} y garantizan el mismo binario en las tres
configuraciones. La imagen de JetPack 7.2.1 no incluye estos módulos en el
intérprete del sistema, de modo que su instalación es explícita y forma parte
del procedimiento registrado. Se usa la compilación de OpenCV de Ubuntu —sin
aceleración por GPU— porque el preprocesamiento y el postprocesamiento se
ejecutan en la CPU en las tres configuraciones; la inferencia, que es la
etapa intensiva en cómputo, se ejecuta en la GPU mediante el runtime que se
fija en la \Cref{sec:pipeline}.

La implementación de middleware es Fast DDS, seleccionada a través de
\texttt{rmw\_fastrtps\_cpp}, que es la implementación por omisión de ROS~2
Jazzy. Se selecciona de forma explícita, en lugar de depender del valor por
omisión de la distribución, para que el middleware efectivo quede registrado
en el entorno y no pueda cambiar de manera inadvertida. La elección se fija
como condición y no como variable porque la implementación de DDS altera de
manera medible la latencia y el uso de recursos \cite{maruyama2016,
kronauer2021, zhang2024}: mientras esa configuración no varíe, las diferencias
observadas entre A, B y C admiten atribución al particionamiento.

No se usa archivo de perfil XML ni se modifican los parámetros de transporte
de Fast DDS: la comunicación entre nodos emplea la configuración por omisión
del middleware que, en un mismo anfitrión, combina el transporte por memoria
compartida con UDP sobre la interfaz de retorno local. Tampoco se habilita la
comunicación intraproceso, porque las tres configuraciones ejecutan los nodos
como procesos separados del sistema operativo. La condición medida corresponde
así al despliegue que resulta de adoptar ROS~2 sin optimizaciones de
transporte, que es precisamente el objeto de estudio.

\subsection{Calidad de servicio}
\label{sec:qos}

Los experimentos usan un único perfil de QoS, idéntico en todos los tópicos y
en las tres configuraciones. El perfil se declara de forma explícita en el
código de los nodos, en lugar de heredarse de los perfiles predefinidos de la
biblioteca cliente, para que sus valores no dependan de la versión instalada y
queden registrados en el artefacto versionado:

\begin{verbatim}
from rclpy.qos import (
    DurabilityPolicy, HistoryPolicy, QoSProfile, ReliabilityPolicy)

PERFIL = QoSProfile(reliability=ReliabilityPolicy.RELIABLE,
                    durability=DurabilityPolicy.VOLATILE,
                    history=HistoryPolicy.KEEP_LAST,
                    depth=1)
\end{verbatim}

La \Cref{tab:qos-policies} resume las políticas fijadas y la alternativa que
se descarta en cada caso.

\begin{table}[htbp]
  \centering
  \caption[Perfil de QoS]{Perfil de QoS fijo. El perfil se aplica a todos los
  tópicos y a las tres configuraciones experimentales.}
  \label{tab:qos-policies}
  \begin{tabular}{@{}L{0.24\linewidth}L{0.20\linewidth}L{0.50\linewidth}@{}}
    \toprule
    \textbf{Política} & \textbf{Valor fijo} & \textbf{Alternativa descartada} \\
    \midrule
    Fiabilidad (\texttt{reliability}) & \texttt{RELIABLE} &
      \texttt{BEST\_EFFORT}: permite al middleware descartar muestras cuando
      el suscriptor se atrasa \\
    Durabilidad (\texttt{durability}) & \texttt{VOLATILE} &
      \texttt{TRANSIENT\_LOCAL}: retiene la última muestra para suscriptores
      que se conectan tarde \\
    Historial y profundidad (\texttt{history}, \texttt{depth}) & \texttt{KEEP\_LAST}, 1 & \texttt{KEEP\_ALL}, o
      una profundidad mayor que uno \\
    Plazo, vigencia y vivacidad (\texttt{deadline}, \texttt{lifespan},
      \texttt{liveliness}) & valores por omisión & declarar plazos,
      vigencias o vivacidades \\
    \bottomrule
  \end{tabular}
\end{table}

Se fija la fiabilidad en \texttt{RELIABLE}, que obliga al middleware a
reintentar la entrega hasta que el suscriptor recibe cada muestra. La
implementación monolítica procesa todos los fotogramas de la entrada, de modo
que la comparación solo es válida si las implementaciones de ROS~2 procesan
esa misma secuencia; con \texttt{RELIABLE} ninguna etapa pierde entradas por
decisión del middleware. La alternativa \texttt{BEST\_EFFORT} no ofrece esa
garantía: cuando una etapa se atrasa, el middleware puede descartar
fotogramas sin notificarlo, de modo que la tasa de procesamiento y la latencia
medidas dejarían de describir el costo arquitectónico para describir también
cuántos fotogramas se descartaron. La opción elegida tiene un costo conocido
—tráfico de confirmaciones y eventuales retenciones del productor cuando el
historial está lleno— que se asume como parte de la condición medida, ya que
la configuración de QoS se refleja en la latencia \cite{fernandez2020,
kronauer2021}.

Se fija la durabilidad en \texttt{VOLATILE}: el suscriptor recibe solo las
muestras publicadas después de emparejar con el publicador. La cadena
transporta un flujo continuo de fotogramas y cada uno solo es útil mientras la
etapa siguiente puede procesarlo, de modo que carece de sentido entregar
muestras anteriores. La alternativa \texttt{TRANSIENT\_LOCAL} mantendría la
última muestra disponible para los suscriptores que se conectan tarde; en una
cadena que arranca por etapas, eso entregaría un fotograma antiguo a la
primera etapa que se incorpora, con lo que la primera medición de cada corrida
incluiría una latencia que no proviene del procesamiento, y obligaría al
publicador a mantener un almacén de muestras, que consume memoria en un
dispositivo de 8~GB.

Se fija el historial en \texttt{KEEP\_LAST} con profundidad 1, de modo que
cada extremo retiene el fotograma en vuelo y nada más: cada etapa procesa un
fotograma a la vez, igual que la implementación monolítica secuencial.
\texttt{KEEP\_ALL} dejaría la cola sin cota, con lo que la memoria crecería
durante la corrida y las latencias incluirían el tiempo de espera en cola, que
es un efecto del tamaño de la cola y no de la arquitectura. Una profundidad
mayor que uno permitiría que dos o más fotogramas estuvieran a la vez en
etapas distintas, es decir, introduciría solapamiento entre etapas: la tasa de
procesamiento aumentaría entonces por la concurrencia añadida y no por el
particionamiento, y la hipótesis H4 —más fronteras implican más sobrecosto—
quedaría confundida, porque el número de fronteras y el grado de solapamiento
cambiarían al mismo tiempo. La profundidad 1 acota además la memoria ocupada
por las imágenes en vuelo: a resolución completa cada muestra ocupa varios
megabytes, y el costo de serialización y de copia crece con el tamaño de la
carga útil \cite{kronauer2021, luo2025}.

Las políticas de plazo, vigencia y vivacidad se dejan en sus valores por
omisión —plazo y vigencia infinitos, vivacidad automática— y se declara de
forma explícita que no se usan. Declarar un plazo o una vivacidad añadiría
vencimientos que compiten con los callbacks de datos dentro del ejecutor y
alterarían el patrón de ejecución que se busca medir; fijar una vigencia, en
cambio, descartaría muestras por antigüedad, que es un descarte por tiempo
con los mismos inconvenientes señalados para \texttt{BEST\_EFFORT}.

El perfil es el mismo en los tópicos de imagen y en el de detecciones, porque
la política de QoS no se compara en este trabajo: una sola configuración evita
introducir un segundo factor experimental. Con esta combinación ninguna
política descarta fotogramas por diseño; si durante los experimentos piloto se
observan pérdidas —por ejemplo, por contrapresión sostenida—, la política de
descarte efectiva se documenta en el procedimiento de la
\Cref{sec:procedure} antes de la campaña final.

La configuración de cómputo se fija junto con el entorno. Los seis núcleos de
la CPU permanecen activos, sin afinidad de procesos ni grupos de control, y el
escalado dinámico de frecuencia se desactiva en los dos dominios de cómputo: el
gobernador de la CPU se fija en \texttt{performance}, que mantiene cada núcleo
en la frecuencia máxima permitida por el modo de energía fijado en la
\Cref{sec:platform}, y el gobernador del dominio de la GPU se fija en el mismo
valor a través de su entrada de \texttt{devfreq}. Ambas condiciones se
verifican antes de cada corrida leyendo el gobernador de cada núcleo y el del
dominio de la GPU, y la frecuencia efectiva de CPU, GPU y memoria se registra
en cada muestra (\Cref{sec:metric-measurement}).

No se deja el escalado por omisión, aunque sea la condición más próxima a un
despliegue sin intervención, porque el escalado dinámico no reacciona a la
arquitectura de la aplicación sino a la carga instantánea de cada dominio, y
las tres configuraciones no imponen la misma carga de CPU: B y C ejecutan más
procesos, más callbacks y más serialización que A. Con el escalado activo, la
frecuencia efectiva de la CPU depende de esa carga, de modo que la
configuración monolítica y las de ROS~2 quedan medidas en regímenes de
frecuencia distintos y la diferencia observada deja de ser atribuible solo al
particionamiento. La variabilidad del escalado no sería entonces ruido que se
promedia, sino un sesgo correlacionado con la variable independiente y
orientado en la dirección de las hipótesis H1 y H3. Fijar la frecuencia en los
dominios de cómputo iguala el régimen en las tres configuraciones y reduce la
fluctuación entre repeticiones, que es la razón por la que la literatura de
medición de ROS~2 neutraliza el escalado de frecuencia del sistema operativo
antes de medir latencia \cite{kronauer2021}. No se recurre a
\texttt{jetson\_clocks}: esa herramienta lleva también el dominio de memoria a
su frecuencia máxima y deja todos los relojes fijos en el extremo del
presupuesto de potencia, condición que se aparta de la establecida en la
\Cref{sec:platform}.

El dominio de memoria conserva su escalado por omisión. Su demanda está
dominada por los mismos tensores de inferencia y los mismos fotogramas en las
tres configuraciones, de modo que su respuesta no se separa sistemáticamente
entre ellas, mientras que fijarlo en el máximo elevaría el consumo y acercaría
al equipo a la limitación por potencia. Esa limitación, si ocurre, se detecta
porque la temperatura y las frecuencias efectivas se registran en todas las
corridas. Si en los experimentos piloto aparecen episodios de limitación
térmica o de potencia, o si la dispersión de las frecuencias registradas
resulta comparable a la señal que se busca medir, la política de relojes se
revisa y se documenta antes de la campaña final.

Para los experimentos principales se utiliza un despliegue sin
contenedores. Todas las condiciones anteriores se tratan como condiciones
experimentales controladas y no como variables bajo investigación; cada una
se documenta de modo que los experimentos puedan reproducirse
(\Cref{sec:reproducibility}).

\section{Cadena de procesamiento de visión por computador}
\label{sec:pipeline}

La carga de trabajo es una cadena de procesamiento de detección de objetos
que consiste en las siguientes etapas lógicas:

\begin{enumerate}
  \item \textbf{Adquisición de imágenes:} obtener cada fotograma de
        entrada.
  \item \textbf{Preprocesamiento:} convertir el fotograma al formato
        esperado por el modelo (p. ej., redimensionamiento, conversión de
        color, normalización y disposición del tensor). Esta etapa es
        individualmente significativa: para la detección de objetos, el
        preprocesamiento contribuye una porción sustancial del tiempo de
        ejecución total, especialmente a resoluciones más altas
        \cite{szegedi2026}.
  \item \textbf{Inferencia:} ejecutar el modelo de red neuronal.
  \item \textbf{Postprocesamiento:} convertir la salida sin procesar del
        modelo en resultados a nivel de aplicación (p. ej., cuadros
        delimitadores, filtrado por confianza, supresión de no máximos). Al
        igual que el preprocesamiento, esta etapa no es despreciable para
        la detección de objetos \cite{szegedi2026}.
  \item \textbf{Salida:} producir el resultado final de la cadena de
        procesamiento.
\end{enumerate}

El modelo de red neuronal se selecciona de acuerdo con los criterios de la
\Cref{sec:workload}, sobre la base de evidencia publicada de que la
detección de objetos en tiempo real es factible en la Orin Nano con las
familias estándar de detectores de una etapa \cite{scalcon2024,
alqahtani2026, szegedi2026,tang2026}:
\todo{SPECIFY: final model, weights version, precision, and inference
runtime once selected.} El mismo modelo y la misma configuración de
inferencia se utilizan en cada configuración experimental.

\section{Implementación monolítica}
\label{sec:monolithic}

La implementación de línea base ejecuta la cadena de procesamiento completa
dentro de un único proceso de Python:

\begin{quote}
\texttt{Entrada $\rightarrow$ Preprocesamiento $\rightarrow$ Inferencia
$\rightarrow$ Postprocesamiento $\rightarrow$ Salida}
\end{quote}

Todas las etapas se ejecutan secuencialmente dentro de un proceso. Los datos
pasan entre las etapas como estructuras de datos en memoria, sin
serialización ni transporte entre procesos. Esta implementación sirve como
línea base de rendimiento contra la cual se comparan las implementaciones de
ROS~2. Corresponde al extremo arquitectónico en el que la aplicación
aprovecha plenamente el hecho de que todos los datos de la cadena pueden
permanecer en un único espacio de direcciones.

\section{Implementación en ROS 2}
\label{sec:ros2-implementation}

Las mismas etapas lógicas se implementan como nodos de ROS~2 utilizando
rclpy, cada uno ejecutándose como un proceso separado del sistema operativo
en el mismo anfitrión, conectados en cadena mediante tópicos que
transportan mensajes de imagen:

\begin{quote}
\texttt{Nodo A $\rightarrow$ Nodo B $\rightarrow$ Nodo C $\rightarrow$
Nodo D}
\end{quote}

Cada nodo publica y se suscribe con el perfil de QoS fijo de la
\Cref{sec:qos}. El mapeo de las etapas de procesamiento a nodos preserva la
equivalencia funcional con la implementación monolítica: dada la misma
secuencia de entrada y el mismo modelo, ambas implementaciones producen
salidas equivalentes.
\todo{SPECIFY: the exact stage-to-node mapping for each configuration once
finalized, and how functional equivalence is verified (e.g., comparing
outputs on a fixed input sequence).}

El ejecutor es idéntico en todos los nodos y en las tres configuraciones
experimentales: cada nodo se ejecuta con un \texttt{SingleThreadedExecutor}
de \texttt{rclpy} y una única suscripción, con el grupo de callbacks
mutuamente exclusivo por omisión y sin temporizadores ni servicios.\footnote{%
La biblioteca cliente de Python expone un ejecutor de un solo hilo y uno de
varios hilos, pero no la variante estática de un solo hilo que sí existe en
\texttt{rclcpp} \cite{rosjazzyrclpy}.} La elección del ejecutor no es
accesoria: el ejecutor determina qué callbacks se ejecutan, en qué hilo y en
qué orden, de modo que forma parte inseparable del rendimiento del nodo
\cite{macenski2023, casini2026}.

Se elige el ejecutor de un solo hilo por tres razones. Primero, procesa los
callbacks en el hilo que invoca a \texttt{spin()} y, por lo tanto, un nodo
procesa un fotograma por vez, igual que la implementación monolítica: un
ejecutor multihilo permitiría que dos fotogramas se procesaran en paralelo
dentro del mismo nodo, lo que añadiría concurrencia interna y confundiría el
efecto atribuible al particionamiento. Segundo, el orden de procesamiento es
determinista y verificable fotograma a fotograma, condición necesaria para
comprobar la equivalencia funcional. Tercero, mantiene el mismo número de
hilos por proceso en las tres configuraciones, de modo que la cantidad de
hilos no varíe junto con el número de fronteras de comunicación.

Esta implementación corresponde a la manera ordinaria y recomendada por
omisión de estructurar una aplicación de ROS~2 como nodos separados
\cite{macenski2022}, antes de que se aplique cualquiera de las mitigaciones
descritas en la literatura relacionada---composición, comunicación
intraproceso o transporte por memoria compartida \cite{macenski2023,
kamiya2025}---. Estudiar esta línea base de forma aislada es deliberado:
mide la arquitectura que una adopción directa de ROS~2 produce.

\section{Configuraciones experimentales}
\label{sec:configurations}

Se utiliza un número reducido de configuraciones, elegidas de modo que
representen diferencias arquitectónicas significativas y no divisiones
arbitrarias de la cadena de procesamiento:

\begin{itemize}
  \item \textbf{Configuración A --- Monolítica:} un proceso de Python.
  \item \textbf{Configuración B --- Particionamiento bajo:} dos nodos de
        ROS~2.
  \item \textbf{Configuración C --- Particionamiento mayor:} cuatro nodos
        de ROS~2.
\end{itemize}

\todo{SPECIFY: the exact node boundaries for Configurations B and C, and
the justification that they represent meaningful architectural
differences.}

La matriz experimental principal cruza las tres configuraciones
arquitectónicas con tres resoluciones de entrada. La resolución se incluye
como dimensión secundaria porque ambos términos del compromiso bajo estudio
escalan con ella: el costo de la inferencia crece con la resolución de
entrada \cite{szegedi2026}, y también el costo de comunicación de los
propios mensajes de imagen, dado que los costos de serialización y copia
por mensaje crecen con el tamaño de la carga útil \cite{kronauer2021,
macenski2022, luo2025}. Las resoluciones se seleccionan según las
capacidades de la fuente de entrada y del modelo elegidos:

\begin{itemize}
  \item Resolución 1: \todo{SPECIFY (candidate: 320$\times$240)}
  \item Resolución 2: \todo{SPECIFY (candidate: 640$\times$480)}
  \item Resolución 3: \todo{SPECIFY (candidate: 1280$\times$720)}
\end{itemize}

Esto da como resultado $3 \times 3 = 9$ condiciones experimentales:

% NOTE: requires \usepackage{amssymb} in the preamble for \checkmark.
\begin{table}[htbp]
  \centering
  \caption[Matriz experimental principal.]{Matriz experimental principal.}
  \label{tab:experimental-matrix}
  \begin{tabular}{clcccc}
    \hline
    \textbf{Config.} & \textbf{Arquitectura} & \textbf{Particionamiento} &
    \textbf{Res.\ 1} & \textbf{Res.\ 2} & \textbf{Res.\ 3} \\
    \hline
    A & Proceso único & 1 proceso & \checkmark & \checkmark & \checkmark \\
    B & ROS 2          & 2 nodos   & \checkmark & \checkmark & \checkmark \\
    C & ROS 2          & 4 nodos   & \checkmark & \checkmark & \checkmark \\
    \hline
  \end{tabular}
\end{table}

Cada condición se repite múltiples veces, o se ejecuta por una duración fija
suficientemente larga, para obtener mediciones estables.
\todo{SPECIFY: number of repetitions and run duration per condition.}

\section{Procedimiento de evaluación comparativa de rendimiento}
\label{sec:procedure}

Cada condición experimental sigue el mismo procedimiento:

\begin{enumerate}
  \item \textbf{Inicializar el sistema.} Iniciar la aplicación y todas las
        herramientas de monitoreo.
  \item \textbf{Calentamiento.} Descartar una cantidad inicial de
        fotogramas o un período fijo de calentamiento para evitar medir el
        comportamiento de inicio: los costos de la primera solicitud, como
        la carga del modelo y la asignación de memoria, no reflejan la
        operación en estado estacionario.
        \todo{SPECIFY: warm-up duration or frame count.}
  \item \textbf{Ejecutar la carga de trabajo.} Ejecutar la cadena de
        procesamiento durante una duración fija o una cantidad fija de
        fotogramas.
        \todo{SPECIFY.}
  \item \textbf{Registrar las mediciones.} Recopilar la latencia de extremo
        a extremo, FPS, la utilización de CPU, RAM y GPU, y el consumo de
        energía si se mide.
  \item \textbf{Repetir.} Repetir el experimento lo suficiente para
        caracterizar la variabilidad.
  \item \textbf{Reiniciar.} Devolver el sistema al mismo estado de
        referencia antes de la siguiente condición experimental.
        \todo{SPECIFY: the reset procedure (e.g., stopping all nodes,
        cooling interval).}
\end{enumerate}

A lo largo de las mediciones, el sistema se mantiene en un estado
controlado: se aplica el modo de energía fijo y se minimiza la carga en
segundo plano. Esto sigue la práctica de la literatura de medición de
ROS~2, donde el comportamiento adaptativo del sistema operativo---el
escalado de frecuencia de la CPU y los estados de inactividad--- se
neutraliza explícitamente porque agrega ruido a las mediciones de latencia
\cite{kronauer2021}.

Los experimentos piloto preceden a los experimentos finales con el fin de
identificar problemas de medición, cuellos de botella inesperados,
supuestos incorrectos y errores de configuración. El propósito de las
ejecuciones piloto es fijar la metodología, no generar resultados finales.

\section{Medición de las métricas de rendimiento}
\label{sec:metric-measurement}

\subsection{Latencia de extremo a extremo}

La latencia se mide por fotograma como el tiempo transcurrido entre la
entrada del fotograma a la primera etapa de la cadena y la disponibilidad de
su salida completamente postprocesada. Las marcas de tiempo se toman en
puntos equivalentes en todas las configuraciones---equivalentes a las
fronteras de la cadena como un todo, no de las etapas internas
individuales--- de modo que las mediciones sean comparables entre
arquitecturas. Esto refleja la perspectiva de extremo a extremo de las
evaluaciones de latencia de ROS~2 a nivel de aplicación \cite{betz2023a} y
la metodología de marcas de tiempo en fronteras utilizada para el análisis
de comunicación resuelto por capa \cite{kamiya2025}.
\todo{SPECIFY: timestamping mechanism (e.g., monotonic clocks) and, for
the ROS 2 configurations, how the timestamps of the first and final nodes
are related on the same host.} Estadísticas informadas: media, mediana,
desviación estándar, mínimo, máximo, P95 y P99 cuando haya muestras
suficientes, siguiendo la práctica de informe de percentiles de la literatura
de evaluación comparativa de rendimiento de ROS~2 \cite{kronauer2021,
macenski2022}.

\subsection{Tasa de procesamiento}

La tasa de procesamiento se mide en fotogramas por segundo sobre el
intervalo de medición en estado estacionario, tras eliminar el calentamiento:

\begin{equation}
  \mathrm{FPS} = \frac{\text{fotogramas procesados}}{\text{tiempo
  transcurrido}}
  \label{eq:fps}
\end{equation}

\subsection{Utilización de CPU}

La utilización promedio y máxima de la CPU se mide durante el intervalo de
medición utilizando herramientas de monitoreo del sistema.
\todo{SPECIFY: tool and sampling interval (e.g., \texttt{tegrastats} or
equivalent).}

\subsection{Utilización de memoria}

La utilización de la RAM del sistema se mide durante el intervalo de
medición; el uso de memoria de GPU se registra si se dispone de una medición
confiable.
\todo{SPECIFY: tool.}

\subsection{Utilización de GPU}

La utilización de GPU se registra para determinar si el sobrecosto
arquitectónico cambia la utilización del recurso de cómputo principal.
\todo{SPECIFY: tool.}

\subsection{Consumo de energía (opcional)}

Si se dispone de una medición confiable: potencia promedio, potencia
máxima, energía por fotograma procesado y rendimiento por vatio, siguiendo
la práctica de métricas de potencia de los estudios de caracterización de la
Orin Nano \cite{tang2026}.
\todo{SPECIFY: measurement method, or state that power measurement was
excluded. Power measurement is an optional extension and must not
threaten completion of the thesis.}

\section{Análisis de datos}
\label{sec:analysis}

Para cada condición experimental, el análisis procede de la siguiente
manera:

\begin{enumerate}
  \item Eliminar las mediciones de inicio y de calentamiento.
  \item Calcular estadísticas descriptivas de la latencia, la tasa de
        procesamiento y la utilización de recursos.
  \item Comparar las distribuciones de latencia entre configuraciones.
  \item Comparar la tasa de procesamiento y la utilización de recursos
        entre configuraciones.
  \item Calcular las diferencias porcentuales entre configuraciones,
        incluido el sobrecosto relativo de cada configuración de ROS~2 con
        respecto a la línea base monolítica:
        \begin{equation}
          \mathrm{sobrecosto}\ (\%) = \frac{L_{\mathrm{ROS}} -
          L_{\mathrm{mono}}}{L_{\mathrm{mono}}} \times 100
          \label{eq:overhead}
        \end{equation}
        donde $L$ denota la métrica bajo comparación (p. ej., la latencia
        media de extremo a extremo).
  \item Investigar si el aumento del particionamiento de nodos produce una
        tendencia consistente a través de las resoluciones.
\end{enumerate}

Donde corresponda, las pruebas estadísticas se seleccionan tras examinar los
datos recopilados y su distribución.
\todo{SPECIFY: tests to be applied once the data distribution is known.}

El análisis evalúa las siguientes hipótesis, enunciadas aquí como
predicciones y no como expectativas que deban confirmarse; los resultados
que contradigan una o más hipótesis siguen siendo válidos y se analizarán
como tales:

% NOTE: The thesis plan lists the hypotheses in the proposal but does not
% assign them a dedicated section in the chapter structure; they are
% placed here because the analysis evaluates them. Move if a better
% location is chosen.
\begin{description}
  \item[H1 (Latencia):] Las implementaciones de ROS~2 exhibirán una mayor
        latencia de extremo a extremo que la implementación monolítica
        debido al sobrecosto de comunicación y procesamiento introducido
        entre los nodos.
  \item[H2 (Tasa de procesamiento):] La implementación monolítica
        alcanzará una tasa de procesamiento igual o mayor que las
        implementaciones de ROS~2 bajo condiciones computacionales
        equivalentes.
  \item[H3 (Utilización de recursos):] Las implementaciones de ROS~2
        requerirán recursos adicionales de CPU y memoria en comparación
        con la implementación monolítica.
  \item[H4 (Particionamiento):] El aumento del número de fronteras de
        comunicación de ROS~2 incrementará el sobrecosto de comunicación
        y, por lo tanto, la latencia de extremo a extremo.
\end{description}

Estas predicciones se fundamentan en la literatura relacionada: el
sobrecosto de comunicación entre procesos se ha medido repetidamente
\cite{kronauer2021, macenski2022, kamiya2025}, y las mitigaciones que
reducen la cantidad de fronteras de proceso recuperan recursos sustanciales
\cite{macenski2023}. Las hipótesis requieren, sin embargo, verificación
experimental para esta carga de trabajo, plataforma y biblioteca cliente,
dado que la magnitud del sobrecosto en relación con el cómputo dominante de
la inferencia es exactamente lo que los estudios previos no aíslan.

\section{Reproducibilidad}
\label{sec:reproducibility}

La siguiente información se documenta de modo que los experimentos puedan
reproducirse, en línea con el llamado de la comunidad a una evaluación
comparativa de rendimiento robótica estandarizada y reproducible
\cite{mayoral-vilches2024}:

\begin{itemize}
  \item Modelo y configuración de hardware
  \item Versión de JetPack
  \item Distribución de Linux
  \item Distribución de ROS~2
  \item Implementación de DDS y versión
  \item Configuración de transporte de DDS
  \item Versión de Python
  \item Versiones de paquetes
  \item Versión del modelo y pesos
  \item Configuración y precisión del modelo
  \item Resolución de entrada por experimento
  \item Configuración de QoS de ROS~2
  \item Tipo de ejecutor y grupos de callbacks
  \item Modo de energía
  \item Configuración de CPU/GPU, gobernador de frecuencia y frecuencias
        efectivas registradas
  \item Condiciones ambientales y del entorno, donde corresponda
  \item Duración del benchmark
  \item Cantidad de repeticiones
  \item Procedimiento de calentamiento
\end{itemize}

Registrar la configuración completa de software y hardware no es meramente
administrativo: las evaluaciones de latencia de sistemas ROS~2 han mostrado
que los resultados dependen de la configuración de software y hardware de
maneras no obvias \cite{betz2023a, kronauer2021}. El código fuente y la
configuración del benchmark están bajo control de versiones. Cualquier
desviación de la metodología planificada que se observe durante los
experimentos finales se registra y se informa.
